\section{A random walk with drift}
\subsection*{(a) Factorization of the conditional density}
With $x_0=0$ fixed, the conditional-density identity gives
\[
f(x_1,\ldots,x_T\mid x_0)=
f(x_T\mid x_{T-1},\ldots,x_0)
f(x_1,\ldots,x_{T-1}\mid x_0).
\]
Iterating proves, for every integer $T\geq1$,
\[
\boxed{f(x_1,\ldots,x_T\mid x_0)
=\prod_{t=1}^{T}f(x_t\mid x_{t-1},\ldots,x_0).}
\]
Independent innovations also give the Markov property:
$x_t\mid(x_{t-1},\ldots,x_0)\sim N(\delta+x_{t-1},\sigma^2)$.

\subsection*{(b) Maximum likelihood estimation}
Let $d_t=x_t-x_{t-1}$. Conditional on $x_0=0$, these increments are
independent $N(\delta,\sigma^2)$ variables, so
\[
\ell(\delta,\sigma^2)
=-\frac T2\log(2\pi)-\frac T2\log\sigma^2
-\frac1{2\sigma^2}\sum_{t=1}^T(d_t-\delta)^2.
\]
The first-order condition for $\delta$ gives
\[
\boxed{\widehat\delta_{\mathrm{ML}}
=\frac1T\sum_{t=1}^Td_t=\frac{x_T}{T}.}
\]
Substituting into the likelihood and differentiating with respect to
$\sigma^2$ gives
\[
\boxed{\widehat{\sigma^2}_{\mathrm{ML}}
=\frac1T\sum_{t=1}^T(d_t-\widehat\delta_{\mathrm{ML}})^2.}
\]
The divisor is $T$, not $T-1$. For $T>1$ these are the interior
maximizers almost surely. For $T=1$, the formal boundary variance
estimate is zero; if $\sigma^2>0$ is required, the likelihood is
unbounded as $\delta=x_1$ and $\sigma^2\downarrow0$, so no interior
variance MLE exists.

\subsection*{(c) Information bound when $\sigma^2=1$}
The score and Fisher information for $\delta$ are
\[
\frac{\partial\ell}{\partial\delta}=\sum_{t=1}^T(d_t-\delta),
\qquad
\mathcal I_T(\delta)
=-E\left[\frac{\partial^2\ell}{\partial\delta^2}\right]=T.
\]
The Cram\'er--Rao bound for an unbiased estimator is therefore $1/T$.
Moreover,
$\widehat\delta_{\mathrm{ML}}\sim N(\delta,1/T)$,
so the estimator in (b) is unbiased and attains the bound.

\clearpage
\subsection*{(d) S\&P 500 monthly index data}
The workbook contains 702 month-end index levels, January 1960--June 2018.
Set $x_0=0$ in January 1960 and
$x_t=\log(SP_t/SP_0)$. This produces $T=701$ increments:
\[
d_t=x_t-x_{t-1}=\log(SP_t/SP_{t-1}).
\]
Applying the formulas in (b),
\[
\boxed{\begin{aligned}\widehat\delta_{\mathrm{ML}}&=0.0055483947,\\
\widehat{\sigma^2}_{\mathrm{ML}}&=0.0017791816,\\
\sqrt{\widehat{\sigma^2}_{\mathrm{ML}}}&=0.0421803461.\end{aligned}}
\]
The fitted average monthly log price growth is approximately $0.555\%$,
and the monthly innovation standard deviation is approximately $4.218\%$.
The corresponding fitted geometric annual price-growth rate is
$\exp(12\widehat\delta)-1\approx6.885\%$.
These are price-index results: dividends are excluded.
They describe the fitted constant-drift Gaussian model, without
establishing that its independence or normality assumptions hold.


\subsection*{(e) Regression on a random walk}
Conditional on the entire regressor vector $X=(x_1,\ldots,x_T)'$,
independence of the errors and the random walk gives
\[
E[\varepsilon\mid X]=0,\qquad
\operatorname{Var}(\varepsilon\mid X)=\gamma^2I_T.
\]
With an intercept, the design matrix has full column rank almost surely
for $T\geq2$, provided $\sigma^2>0$. Thus the model satisfies the
\emph{conditional} Gauss--Markov assumptions. Serial dependence of
the regressors does not invalidate these conditions.

The pairs $(y_t,x_t)$ are not i.i.d. Indeed,
\[
E[x_t]=t\delta,\qquad
\operatorname{Var}(x_t)=t\sigma^2,\qquad
\operatorname{Cov}(x_s,x_t)=\sigma^2\min(s,t).
\]
The distribution of $x_t$ changes with $t$, and the regressors are
serially dependent.

For a two-sided test of $H_0:\beta=1$ with unknown $\gamma^2$,
fit the model by OLS and compute
\[
t_{\mathrm{obs}}=
\frac{\widehat\beta-1}{s/\sqrt{\sum_{t=1}^T(x_t-\bar x)^2}},
\qquad
s^2=\frac{\sum_{t=1}^T
(y_t-\widehat\alpha-\widehat\beta x_t)^2}{T-2}.
\]
Conditional normality implies the exact null distribution
$t_{\mathrm{obs}}\mid X\sim t_{T-2}$. Since this distribution does not
depend on $X$, it also holds unconditionally.
For $T=3$, use $t_1$; at the two-sided $5\%$ level reject if
$|t_{\mathrm{obs}}|>12.706$.
